\section{Background: What And Where Do Diffusion Models Learn?} \label{sec:background}
In this section, we review the formulation of diffusion models along with the commonly held interpretation, which we refer to as \emph{global underfitting}. This sets the stage for our contrasting viewpoint--\emph{selective underfitting}--introduced in the following sections.

Diffusion models\footnote{Throughout this work, we use ``diffusion'' as an umbrella term for both diffusion and flow matching~\citep{lipman2022flow,liu2022flow}, as they are equivalent~\citep{gao2025diffusion}. All discussions apply to both.} define a stochastic forward process that transforms a data distribution into a Gaussian distribution over timesteps $t \in [0, 1]$: $\rvz_t = \alpha_t \rvx + \sigma_t \rveps$ where $\rvx$ is a clean data and $\rvepsilon \sim \gN(\vzero, \mI)$ for schedules $\alpha_t,\sigma_t \colon [0,1] \to \mathbb{R}$. The models are trained through denoising score matching \venueifelse{(DSM, ~\citet{vincent2011connection})}{(DSM)~\citep{vincent2011connection}} to learn the score function--the gradient of the log density--which is then used at inference to iteratively denoise Gaussian noise into a sample from the data distribution.

\begin{figure}
    \centering
    % \includegraphics[width=\textwidth,height=0.35\textwidth]{example-image-duck}
    \begin{overpic}[width=0.64\textwidth]{figures/extrapolation/extrapolation_illust}
    % \put(4, 49){\footnotesize\xgray{(a)}}
    \end{overpic}
    \hfill
    \begin{overpic}[height=0.34\textwidth,trim={1cm 0.15cm 0 0}]{figures/extrapolation/extrapolation}
    \end{overpic}
    \venueifelse{\vskip -0.5\baselineskip}{\vskip -0.25\baselineskip}
    \caption{\textbf{Selective underfitting.} (a) Training distribution concentrates in a \coloredul{xblue}{small region} of data space, where the model learns the \xxgreen{empirical score function $\rvs_\star$}. At inference, \coloredul{xred}{sampling trajectories} go beyond this region, where predictions are not directly \xxgreen{\emph{\textbf{supervised}}} and must be \xsienna{\emph{\textbf{extrapolated}}}. (b) The plotted distances reflect how far \coloredul{xred}{sampling trajectories} are from the \coloredul{xblue}{supervision region}.}
    \label{fig:extrapolation}
    \phantomsubcaption\label{fig:extrapolation-a}
    \phantomsubcaption\label{fig:extrapolation-b}
    % \neuripsonly{\vskip -0.5\baselineskip}
    \vskip -\baselineskip
\end{figure}

\textbf{Learning Problem.} Given training data $\train = \left\{\rvx^{(i)}\right\}_{i=1}^N$, the DSM objective at timestep $t$ is:
\begin{equation}\label{eq:dsm}
\Ls_{\text{DSM}} (t) := \E_{\rvx^{(i)} \sim \mathrm{Unif}(\train), \rveps \sim \gN(\vzero, \mI)} \left\| \rvs_\rvtheta (\rvz_t, t) + \rveps / \sigma_t \right\|^2
\end{equation}
The DSM nature of diffusion model training can be decomposed into two components: \xxgreen{(i) \emph{what function}} does the model aim to learn? (i.e., what is the minimizer of $\Ls_{\text{DSM}}$?), and \xblue{(ii) \emph{where in the data space}} does the model learn during training? (i.e., what is the distribution of $\rvz_t$?). For (i), there exists an analytic optimal solution for \Cref{eq:dsm}:
\begin{equation} \label{eq;empirical_score}
\rvs_\star (\rvz_t, t) = \frac{1}{\sigma_t^2} \left[ -\rvz_t + \alpha_t \sum_{i=1}^{N} \operatorname{softmax}\left(-\frac{\|\rvz_t - \alpha_t \rvx^{(i)}\|^2}{2 \sigma_t^2}\right) \rvx^{(i)} \right].
\end{equation}
Since this is the score function of the empirical training data, we refer to $\rvs_\star$ as the \emph{empirical score function}, often also called the analytic score function. For (ii), a model is trained on noisy versions of the training data, which can be written as a mixture of Gaussians: $\hat{p}_t(\rvz_t) = \frac{1}{N}\sum_{i=1}^{N}\gN(\rvz_t;\alpha_t \rvx^{(i)}, \sigma_t^2 \mI)$. Given (i) and (ii), \Cref{eq:dsm} is equivalent to:
\begin{equation} \label{eq:dsm-rewritten}
\Ls_{\text{DSM}} (t) = \E_{\xblue{\rvz_t \sim \hat{p}_t}} \left\|\rvs_\rvtheta (\rvz_t, t) - \xxgreen{\rvs_\star (\rvz_t, t)}\right\|^2 + (\text{constant}).
\end{equation}
In principle, \emph{$\hat{p}_t$ has support over the entire data space} for any $t>0$, since it is Gaussian. This leads to the following common understanding of the learning problem:
\begin{perspective}[Common] \label{perspective:common}
    Diffusion models are supervised to \xxgreen{(i) approximate the empirical score function $\rvs_\star$,} \xblue{(ii) over the entire data space.}
\end{perspective}
\newcommand{\common}{\Cref{perspective:common}\xspace}

\common, which we call \textbf{global underfitting}, has been widely adopted to study diffusion models, for example, in the context of their generalization~\citep{kamb2024analytic,kadkhodaie2023generalization,scarvelis2023closed,yi2023generalization,wang2024on} and memorization~\citep{niedoba2024towards}. These works primarily focus on (i), analyzing \emph{how} a neural network globally underfits the empirical score function. In contrast, we argue that (ii)--\emph{where} learning occurs--is the critical component, as underfitting occurs \emph{selectively} on certain regions of the data space.